\section{Finite-q Wannier kernel and gauge covariance}
\label{sec:finite-q}

Using the Fourier conventions in \cref{eq:atomic-bloch}, the DFPT vertex transfers an electron from $\kk$ to $\kk+\qq$, while the spin vertex at $-\qq$ closes the loop. Define the complex retarded loop
\begin{equation}
\begin{split}
\bA^R_{\ell a,\kappa\mu}(\qq)
=&\frac{1}{N_k}\int\dd\varepsilon\,f(\varepsilon)
\sum_{\kk}\tr\Big[
\bT_{\ell a}(-\qq;\kk,\kk+\qq)
G^R_{\kk+\qq}(\varepsilon)\\
&\hspace{35mm}\times
g_{\kappa\mu}(\kk,\qq)
G^R_{\kk}(\varepsilon)\Big].
\end{split}
\label{eq:retarded-loop}
\end{equation}
This is the q-resolved Wannier realization proposed here. The phonon-like effective field of Ref.~\cite{Mankovsky2023SLC} motivates the finite-q object, but the explicit momentum-space DFPT loop is a project extension.

A real-space Hermitian response need not have a real Fourier coefficient at a generic q. The physical kernel is obtained from the retarded/advanced q pair,
\begin{equation}
K^{\rm bub}_{\ell a,\kappa\mu}(\qq)
=-\frac{1}{2\pi\ii}
\left[
\bA^R_{\ell a,\kappa\mu}(\qq)
-\bA^R_{\ell a,\kappa\mu}(-\qq)^*
\right].
\label{eq:qpair-kernel}
\end{equation}
It follows that
\begin{equation}
K(-\qq)=K(\qq)^*.
\label{eq:q-reality}
\end{equation}
At a self-inverse q point, \cref{eq:qpair-kernel} reduces to the familiar $-\operatorname{Im}\bA^R/\pi$. Applying an elementwise imaginary part to one generic-q loop would incorrectly discard the phase of the Fourier coefficient.

If $\kk+\qq=\overline{\kk+\qq}+\bm G$, the atomic-position gauge uses
\begin{equation}
D_{\bm G}=\operatorname{diag}_n
 e^{\ii\bm G\cdot\bm\tau_n}\otimes I_{\rm spin},
\end{equation}
and
\begin{equation}
G(\kk+\qq)=D_{\bm G}^\dagger
G(\overline{\kk+\qq})D_{\bm G}.
\label{eq:wrap-green}
\end{equation}
The final-state index of $g(\kk,\qq)$ transforms consistently. After conversion, Hermitian lattice perturbations obey
\begin{equation}
g(\kk,\qq)^\dagger=g(\kk+\qq,-\qq),
\label{eq:g-hermiticity}
\end{equation}
with the appropriate wrapped indices.
